\documentclass[12pt]{article}
% Urdu (right to left, Nastaliq script) with English and math, using
% polyglossia and XeLaTeX as most Overleaf Urdu templates do. Awami Nastaliq
% is the Nastaliq font TeX Live ships; it is rendered with Graphite.
\usepackage{amsmath}
\usepackage{polyglossia}
\setdefaultlanguage{urdu}
\setotherlanguage{english}
\setmainfont{Libertinus Serif}
\newfontfamily\urdufont[Script=Arabic,Renderer=Graphite,Scale=0.9]{Awami Nastaliq}
\newfontfamily\urdufontsf[Script=Arabic,Scale=1.0]{Amiri}
\begin{document}

\section{تعارف}
یہ ایک آزمائشی دستاویز ہے جو اردو کو \textenglish{English} کے ساتھ ملاتی ہے۔
پاکستان کا دارالحکومت اسلام آباد ہے، اور سب سے بڑا شہر کراچی ہے۔

\section{ریاضی}
آئن سٹائن کی مساوات یہ ہے:
\begin{equation}
  E = mc^2
\end{equation}
اور سال ۲۰۲۶ میں ۴۲ مخزن آزمائے گئے۔

\begin{english}
\section{English paragraph}
This paragraph is left to right inside a right-to-left document, with
\texturdu{اردو} words inside it.
\end{english}

\begin{itemize}
  \item پہلا نکتہ
  \item دوسرا نکتہ: {\urdufontsf نسخ خط میں متن}
\end{itemize}

\end{document}
